Figure 1 shows the light curve of the eclipsing binary V1107 Cas, classified
as a W Ursae Majoris type. Table 1 contains a list of observed minima of the
light variation. The columns contain: the number of the minimum, the date on
which the minimum was observed, the heliocentric time of minimum expressed in
Julian days and an error (in fractions of a day).

Using these data:
\begin{description}
\item[(a)] Determine an initial period of V1107 Cas, assuming that the
  period of the star is constant during the interval of observations. Assume
  that observations during one night are continuous. Duration of the transit
  is negligible.
\item[(b)] Make what is known as an (O$-$C) diagram (for ``observed $-$
  calculated'') of the times of minima, as follows: on the $x$-axis put the
  number of periods elapsed (the ``epoch'') since a chosen initial moment
  $M_o$; on the $y$-axis the difference between the observed moment of
  minimum $M_{\mathrm{obs}}$ and the moment of minimum calculated using the
  formula (``ephemeris''):
  \[ M_{\mathrm{calc}} = M_o + P \times E \]
  where $E$, the epoch, is exactly an integer or half-integer, and $P$ is
  the period in days.
\item[(c)] Using this (O$-$C) diagram, improve the determination of the
  initial moment $M_o$ and the period $P$, and estimate the errors in their
  values.
\item[(d)] Calculate the predicted times of minima of V1107 Cas given in
  heliocentric JD occurring between 19h, 1 September 2011 UT and 02h,
  2 September 2011 UT.
\end{description}

\begin{center}
\begin{tabular}{cccc}
No. & Date of minimum (UT) & Time of minimum (Heliocentric JD) & Error \\
1 & 22 December 2006 & 2\,454\,092.4111 & 0.0004 \\
2 & 23 December 2006 & 2\,454\,092.5478 & 0.0002 \\
3 & 23 September 2007 & 2\,454\,367.3284 & 0.0005 \\
4 & 23 September 2007 & 2\,454\,367.4656 & 0.0005 \\
5 & 15 October 2007 & 2\,454\,388.5175 & 0.0009 \\
6 & 15 October 2007 & 2\,454\,388.6539 & 0.0011 \\
7 & 26 August 2008 & 2\,454\,704.8561 & 0.0002 \\
8 & 5 November 2008 & 2\,454\,776.4901 & 0.0007 \\
9 & 3 January 2009 & 2\,454\,835.2734 & 0.0007 \\
10 & 15 January 2009 & 2\,454\,847.3039 & 0.0004 \\
11 & 15 January 2009 & 2\,454\,847.4412 & 0.0001 \\
12 & 16 January 2009 & 2\,454\,847.5771 & 0.0004 \\
\end{tabular}

\medskip
Table 1: Observed times of minima of V1107 Cassiopeae
\end{center}
